\begin{abstract}
Common knowledge is crucial for safe group coordination. In its absence, %When common knowledge cannot be obtained, 
humans must %are forced to 
rely on shared knowledge, %only.
which is inherently limited in depth and therefore prone to coordination failures,
%humans rely on their ability to identify shared knowledge to distinguish who knows what.
%When common knowledge can not be obtained, humans routinely use their ability to attribute beliefs and intentions \tobo{What does intentions have to do with it? Next sentence also says "*such* shared knowledge attributions", which in principle is neither about beliefs nor intentions.} \tsn {Maybe I misunderstand, but intentions are related due to the fact that coordination relies on being able to correctly attribute (predict) the intentions of others. But I find another issue: It might imply that people *only* attribute beliefs to others when common knowledge cannot be obtained. But every social action involves belief attribution (or attributing some mental state to others). So, could word it as something like "When common knowledge cannot be obtained, humans rely on their ability to identify shared knowledge to distinguish between who knows what" }  in order to infer what is known. 
%But shared knowledge is limited in depth and therefore prone to coordination failures. %
because any finite-order knowledge attribution allows for an even higher order attribution that may change what is known by whom. 
\modified{In three separate experiments involving 134 participants, we investigate the extent to which humans
can differentiate between common knowledge and $n$th-order shared knowledge.} We %have 
designed a two-person coordination game with imperfect information to simplify the %that is able to cast the 
recursive game structure  and higher-order uncertainties into a relatable everyday scenario. 
In this game, coordination %to achieve 
for the highest payoff requires a specific fact to be common knowledge between players. 
However, this fact cannot become common knowledge %is impossible to achieve 
in the game. The fact can at most be $n$th-order shared knowledge for some $n$.
%only nth-order shared knowledge can be achieved for varying values of n.
Our findings reveal that even at quite shallow depths of shared knowledge (low values of $n$), players behave as though they possess common knowledge, and %express 
claim similar levels of certainty in their actions, %as if they had. 
despite incurring significant penalties when falsely assuming guaranteed coordination.
We term this phenomenon `the curse of shared knowledge'. It %occurs in spite of participants suffering huge payoff penalties when falsely assuming that coordination, and hence common knowledge, is guaranteed.
  %\tsn{Måske korte lidt ned, eller hive lidt i land? Synes ikke vi konkluderer for stærkt her, udover måske "express certainty", som måske altid er lidt svært at rigtigt vide.} \add{
 % We conjecture that this is 
arises either from the players' %participants' 
inability to distinguish between higher-order shared knowledge and common knowledge, or from their implicit assumption that their co-player cannot make this distinction.
  
%express to be completely certain about being coordinated and achieving the highest payoff,  
% express to be as certain that they are coordinated and will achieve the highest payoff as if they indeed had common knowledge. 
%Despite this, our results show that participants behave as if they had common knowledge, and for sufficiently large values of n, they claim to  
% in self-assessing the certainty of receiving the 
%their actions
%claim their actions to be as certain as if they had common knowledge. 
%It is impossible to achieve common knowl
%
%Our results show that participants have a very hard time accepting the fact that common knowledge is not reducible to shared knowledge.\tobo{Maybe too strong? Do our results really show that? They behave as if having common knowledge, but they might still be able to distinguish common knowledge achieved via public announcements from what they have in this game.} Instead, participants try to coordinate even at quite shallow depths of shared knowledge and in spite of huge payoff penalties. \tsn{With regards to Tobo's comments consider the following alternative: "Our results show that participants try to coordinate when the levels of shared knowledge gets sufficiently high. In other words, our participants behave as if they have common knowledge about facts that are only shared knowledge, and collectively suffer huge payoff penalties as a result."} \re{Yes too strong, i guess. Proposal, taget fra TOBO: "Our results show that participants behave as if they had common knowledge. Participants also claim to be certain about their actions as if they had common knowledge, even when common knowledge is not obtainable."}



%\linenumbers

% Use "Eq" instead of "Equation" for equation citations.
%\section*{Introduction}


% Results and Discussion can be combined.
%\section*{Results}

%PLOS does not support heading levels beyond the 3rd (no 4th level headings)
\end{abstract}  
